\section{Related Work}
\label{sec:related}

\para{Multi-turn evaluation and degradation.}
A growing body of work documents systematic performance drops in multi-turn settings. \citet{laban2025lost} introduce a ``sharding'' framework that decomposes single-turn tasks into multi-turn equivalents and find a 39\% average performance drop across 15 models and 6 tasks. Their \concat control---which presents all information in a single message---recovers 95.1\% of single-turn performance, isolating the degradation to sequential disclosure rather than information content. \citet{myung2026conversational} confirm this pattern with controlled experiments on instruction following, tool selection, and entity extraction, reporting that \gptmini drops from 93\% to 24\% on instruction following across turns. \citet{sirdeshmukh2025multichallenge} find that even the best frontier model scores only 41.4\% on their multi-turn benchmark, with instruction retention being a primary failure mode. Our work extends this line by asking not just \emph{whether} models degrade, but whether the degradation mechanism involves regression toward the base model's distribution.

\para{The superficial alignment hypothesis.}
\citet{lin2024urial} demonstrate that alignment tuning affects only 5--8\% of output tokens, predominantly stylistic tokens such as greetings, hedges, and refusal phrases. They show that KL divergence between base and aligned models decreases monotonically over token position within a response, and that simple in-context learning (3 examples) can replicate most alignment behaviors without any fine-tuning. \citet{li2024ssah} complement this finding at the neuron level, showing that safety alignment localizes to roughly 1.4\% of neurons (Safety Critical Units) that are fragile under fine-tuning---65.7\% of these neurons get reassigned during downstream adaptation. \citet{vergara2026superficial} formalize the superficial alignment hypothesis through Kolmogorov complexity, showing that post-training adds only kilobytes of information atop terabytes of pre-training knowledge. These findings collectively suggest that alignment is a thin veneer over the base distribution, motivating our hypothesis that extended interaction may erode this veneer.

\para{Multi-turn safety and adversarial attacks.}
\citet{li2026multiturn_jailbreak} demonstrate that multi-turn attacks using tree search achieve 97\% success rates against aligned models by gradually shifting conversational context. While their work focuses on adversarial settings, it demonstrates that alignment can be eroded through sequential interaction even when individual turns appear benign. Our work complements this by testing whether alignment degrades \emph{without} adversarial intent---through natural conversation alone.

\para{Alignment training for multi-turn settings.}
\citet{shi2024multiturn_rlhf} show that applying single-turn RLHF objectives to multi-turn settings introduces covariate shift, as the training distribution does not match the sequential nature of real conversations. This training-deployment mismatch provides a mechanistic explanation for why alignment might weaken at later turns: the model has simply not been trained to maintain alignment in extended interaction. Our experiments provide direct empirical evidence for the consequences of this mismatch.
